[ignoring loop detection]
\documentclass[runningheads]{llncs}

\usepackage[T1]{fontenc}
\usepackage[demo]{graphicx} 
\usepackage{booktabs}
\let\vec\relax
\usepackage{amsmath}
\usepackage{amssymb}
\usepackage{multirow}
\usepackage{xcolor}
\usepackage{subcaption}
\usepackage[ruled,vlined,linesnumbered]{algorithm2e}
\usepackage[pdfusetitle=false]{hyperref}
\hypersetup{pdftitle={ArtBench Generative Modeling: A Comparative Study of VAEs, GANs, and Diffusion Models}}

\begin{document}

\title{ArtBench Generative Modeling:\texorpdfstring{\\}{ }A Comparative Study of VAEs, GANs, and Diffusion Models}

\author{Alexandre Rodrigues\inst{1} \and
Duarte Pereira\inst{1}}

\authorrunning{A. Rodrigues et D. Pereira.}

\institute{Faculdade de Ciências e Tecnologia, Universidade de Coimbra, Portugal\\
\email{2022236193@student.uc.pt, xxx@student.uc.pt}}

\maketitle

% ---------------------------------------------------------------
\begin{abstract}
We compare three families of deep generative models --- Variational Autoencoders (VAEs), Deep Convolutional GANs (DCGANs), and Denoising Diffusion Probabilistic Models (DDPMs) --- on the ArtBench-10 dataset ($32{\times}32$, 10 art styles, 50\,000 images). Over 100 independent training runs were conducted via an automated experiment orchestrator, followed by full-dataset retraining of the best configuration per family. Under a unified evaluation protocol (FID and KID, 5\,000 samples, 10 seeds), the VAE achieves FID\,=\,145.85, the DCGAN with Spectral Normalization achieves FID\,=\,28.12\,$\pm$\,0.59, and the Diffusion model with EMA and DDIM sampling achieves FID\,=\,28.97\,$\pm$\,0.38. At $32{\times}32$ resolution, a well-tuned GAN is competitive with diffusion models, both outperforming the VAE by over 117 FID points.
\keywords{Generative Models \and VAEs \and GANs \and Diffusion \and ArtBench \and FID}
\end{abstract}

% ---------------------------------------------------------------
\section{Introduction}\label{sec:intro}

Generative modeling of visual art presents distinctive challenges: artwork images exhibit strong stylistic structure, high intra-class variability, and no privileged semantic object to anchor the distribution. The ArtBench-10 dataset~\cite{liao2022} was specifically designed to benchmark generative models in this domain, covering 10 artistic styles from Baroque to Ukiyo-e at $32{\times}32$ resolution.

This work presents a systematic comparison of three model families on ArtBench-10: VAEs~\cite{kingma2013}, DCGANs~\cite{radford2015}, and DDPMs~\cite{ho2020}. Beyond the mandatory baselines, we explore extensions including Spectral Normalization~\cite{miyato2018}, WGAN-GP~\cite{gulrajani2017}, conditional DCGAN, StyleGAN~\cite{karras2018}, DDIM sampling~\cite{song2020}, EMA, CVAE, and VQ-VAE~\cite{oord2017}. Our goal is to understand the technical trade-offs between stability, sampling speed, and visual fidelity.

A key contribution is the engineering infrastructure: a profile-based configuration system (\texttt{config.py}) with environment variable overrides and a custom orchestrator (\texttt{run\_experiments.py}) that enabled 100+ reproducible experiments across three machines (Apple M1, M4, and NVIDIA RTX 3060) with minimal manual intervention.

% ---------------------------------------------------------------
\section{Problem Statement and Dataset}\label{sec:problem}

The task is \emph{unconditional image generation}: learning a model that can sample new images statistically indistinguishable from real artwork. Unlike classification, there is no ground-truth label; quality must be assessed both visually and through distribution-level metrics. The challenge is threefold: achieving visual fidelity (coherent structure and texture), diversity (avoiding mode collapse), and style plausibility (samples should resemble painted artworks).

ArtBench-10~\cite{liao2022} contains 50\,000 training images (5\,000 per class) and 10\,000 test images across ten artistic styles. All images are $32 \times 32$ RGB, normalized to $[-1, 1]$. During development, a fixed 20\% stratified subset (10\,000 images) is used for hyperparameter search, ensuring reproducibility across machines. The best configuration per family is then retrained on the full dataset for final evaluation.

% ---------------------------------------------------------------
\section{Methodology}\label{sec:methodology}

\subsection{Variational Autoencoder (VAE)}
The ConvVAE trains an encoder--decoder pair jointly. The encoder maps an image to a diagonal Gaussian distribution $(\mu, \sigma^2)$ in latent space ($d=128$). Training minimizes the $\beta$-VAE objective~\cite{higgins2017}: MSE reconstruction loss plus $\beta$ times the KL divergence from the standard normal prior. The architecture has $\approx$1M parameters. We also implemented a \textbf{Conditional VAE} (style-conditioned via learned embeddings) and a \textbf{VQ-VAE} (discrete codebook with 256 entries).

\subsection{Generative Adversarial Network (DCGAN)}
The DCGAN~\cite{radford2015} frames generation as a minimax game between Generator ($G$) and Discriminator ($D$). We use \textbf{Spectral Normalization (SN)}~\cite{miyato2018} to enforce a 1-Lipschitz bound on $D$ weight matrices, preventing discriminator saturation. We also explore \textbf{WGAN-GP}~\cite{gulrajani2017}, \textbf{conditional DCGAN}, and \textbf{StyleGAN}~\cite{karras2018}.

\subsection{Diffusion Model (DDPM)}
The DDPM~\cite{ho2020} reverses a Markov noising process using a pixel-space U-Net ($\approx$10.5M parameters). Inference uses \textbf{DDIM}~\cite{song2020} with 100 deterministic steps, summarized in Algorithm~\ref{alg:ddim}. \textbf{EMA} (decay $=0.9999$) is used for final sampling.

\begin{algorithm}[htbp]
\DontPrintSemicolon
\SetAlgoLined
\KwIn{Trained noise predictor $\varepsilon_\theta$; steps $S=100$; $\eta=0$}
Sample $x_T \sim \mathcal{N}(0,I)$\;
\For{$i = S$ \KwTo $1$}{
  $\hat{\varepsilon} \leftarrow \varepsilon_\theta(x_{\tau_i}, \tau_i)$; $\hat{x}_0 \leftarrow \bigl(x_{\tau_i} - \sqrt{1-\bar{\alpha}_{\tau_i}}\,\hat{\varepsilon}\bigr)/\sqrt{\bar{\alpha}_{\tau_i}}$\;
  $x_{\tau_{i-1}} \leftarrow \sqrt{\bar{\alpha}_{\tau_{i-1}}}\,\hat{x}_0 + \sqrt{1-\bar{\alpha}_{\tau_{i-1}}}\,\hat{\varepsilon}$\;
}
\Return $x_0$\;
\caption{DDIM Sampling (Deterministic)}
\label{alg:ddim}
\end{algorithm}

% ---------------------------------------------------------------
\section{Experiment Orchestration}\label{sec:engineering}

To manage 100+ experiments, we developed a profile-based system:
\noindent\textbf{Config System (\texttt{config.py}).} Defines TEST (validation), DEV (30-50 epochs, 20\% subset, 3 seeds) and PROD (200 epochs, full dataset, 10 seeds) profiles. Hyperparameters are injectable via environment variables, allowing grid searches without code modification.
\noindent\textbf{Orchestrator (\texttt{run\_experiments.py}).} Automates sequential training, triggers evaluation, and saves isolated result directories. This ensured reproducibility across heterogeneous machines.

% ---------------------------------------------------------------
\section{Experimental Setup}\label{sec:setup}

Models are evaluated using FID~\cite{heusel2017} and KID~\cite{binkowski2018} over 5\,000 generated samples. KID uses an unbiased MMD$^2$ estimator (50 subsets of 100 images), making it preferred for smaller evaluation sets. Table~\ref{tab:search_all} summarizes the search space.

\begin{table}[htbp]
\caption{Hyperparameter search space per model family.}\label{tab:search_all}
\centering\small
\begin{tabular}{l|ll|ll|ll}
\toprule
& \multicolumn{2}{c|}{VAE (45+ runs)} & \multicolumn{2}{c|}{DCGAN (35 runs)} & \multicolumn{2}{c}{Diffusion (25+ runs)} \\
Param & Searched & Best & Searched & Best & Searched & Best \\
\midrule
$\beta$ / $C$ / lat & 0--2.0 & \textbf{0.05} & lat: 32--256 & \textbf{100} & $C$: 32--128 & \textbf{96} \\
Latent $d$ & 16--256 & \textbf{128} & ngf: 32--128 & \textbf{128} & $T$: 100--2000 & \textbf{1000} \\
LR & 1e-4--1e-2 & \textbf{2e-3} & 2e-4, 1e-3 & \textbf{2e-4} & 2e-5--1e-3 & \textbf{4e-4} \\
Epochs & 30--200 & \textbf{200} & 50--200 & \textbf{200} & 50--200 & \textbf{200} \\
Key tech & cosine LR & \textbf{---} & SN, cos LR & \textbf{SN} & cos LR, EMA & \textbf{both} \\
\bottomrule
\end{tabular}
\end{table}

% ---------------------------------------------------------------
\section{Results and Ablation Studies}\label{sec:results}

\subsection{VAE Ablation}\label{sec:vae_ablation}

\subsubsection{KL Weight $\beta$ --- The Most Critical Hyperparameter.}
The FID-versus-$\beta$ curve (Table~\ref{tab:vae_beta}) is U-shaped with an optimum at $\beta \in [0.05, 0.1]$. At $\beta=0$, the model is a pure autoencoder with no latent regularity. At $\beta \geq 0.5$, \textbf{posterior collapse} causes blurry outputs. The ArtBench optimum is far below the notebook default of 0.7, reflecting that 10 art styles need relaxed regularization to capture texture.

\begin{table}[htbp]
\caption{VAE $\beta$ sweep ($d=128$, lr$=2{\times}10^{-3}$, DEV).}\label{tab:vae_beta}
\centering\small
\begin{tabular}{ccc|ccc}
\toprule
$\beta$ & FID (30ep) & FID (150ep) & $\beta$ & FID (30ep) & FID (150ep) \\
\midrule
0.0  & 336.8 & --- & 0.5 & 223.2 & --- \\
0.02 & ---   & 148.2 & 0.7 & 241.6 & --- \\
\textbf{0.05} & 205.1 & \textbf{140.2} & 1.0 & 262.3 & --- \\
0.1  & 185.3 & 141.3 & 2.0 & 316.3 & --- \\
\bottomrule
\end{tabular}
\end{table}

\subsubsection{Other VAE Ablations.}
\begin{table}[htbp]
\caption{VAE ablation over latent dim, LR, and epochs (DEV subset).}\label{tab:vae_other}
\centering\small
\begin{tabular}{cc|cc|cc}
\toprule
\multicolumn{2}{c|}{Latent dim ($\beta$=0.7)} & \multicolumn{2}{c|}{Learning rate ($\beta$=0.7)} & \multicolumn{2}{c}{Epochs ($\beta$=0.1, lr=2e-3)} \\
$d$ & FID & LR & FID & Epochs & FID \\
\midrule
16  & 356.3 & 1e-4 & 443.7 & 30  & 169.6 \\
64  & 234.1 & 1e-3 & 241.6 & 100 & 146.1 \\
128 & 241.1 & \textbf{2e-3} & \textbf{207.4} & 150 & \textbf{141.3} \\
256 & 240.9 & 1e-2 & 456.6 & 200 (PROD) & \textbf{145.9} \\
\bottomrule
\end{tabular}
\end{table}
\textbf{Latent dim:} FID saturates at $d \geq 128$. \textbf{LR:} Optimal at 2e-3. \textbf{Cosine LR:} Hurt performance (+6 FID) as the LR fell below $10^{-4}$ too early, stalling learning.

\subsection{DCGAN Ablation}\label{sec:gan_ablation}

\subsubsection{Capacity and Latent Dim.} Smaller latent spaces ($d=32$) help the GAN structure its mapping. ngf=128 beats ngf=64 by 14 FID points. Adam $\beta_1=0.9$ raises FID from 119 to 287, confirming that lower momentum is essential for adversarial stability.

\begin{table}[htbp]
\caption{DCGAN ablations (DEV, 50 epochs).}\label{tab:dcgan_ablation}
\centering\small
\begin{tabular}{lcc|lcc}
\toprule
\multicolumn{3}{c|}{Latent dim (ngf=64)} & \multicolumn{3}{c}{G/D capacity (lat=100)} \\
Lat & FID & KID & ngf/ndf & FID & KID \\
\midrule
\textbf{32} & \textbf{103.4} & 0.068 & 32/32 & 139.3 & 0.113 \\
100 & 118.8 & 0.090 & \textbf{128/128} & \textbf{104.8} & \textbf{0.073} \\
\bottomrule
\end{tabular}
\end{table}

\subsubsection{Spectral Normalization --- The Decisive Stabilizer.}
SN enforces the Lipschitz bound, preventing \textbf{discriminator saturation} (D loss $\to 0$). This enabled stable 200-epoch training (Table~\ref{tab:dcgan_sn}).
\begin{table}[htbp]
\caption{Impact of Spectral Normalization (ngf/ndf=128, lat=100).}\label{tab:dcgan_sn}
\centering\small
\begin{tabular}{lccc|lccc}
\toprule
Config & Epochs & FID & KID & Config & Epochs & FID & KID \\
\midrule
Without SN & 50  & 104.8 & 0.073 & \textbf{With SN} & 100 & \textbf{71.2} & \textbf{0.034} \\
Without SN & 100 & 74.3  & 0.039 & With SN & 200 & \textbf{28.1} & \textbf{0.011} \\
\bottomrule
\end{tabular}
\end{table}

\subsection{Diffusion Ablation}\label{sec:diff_ablation}

\textbf{Cosine LR scheduling} was transformative (Table~\ref{tab:diff_path}), saving 115 FID points. \textbf{EMA} requires a long convergence horizon, achieving the best results only at epoch 200.
\begin{table}[htbp]
\caption{Progressive diffusion improvement ($C=96$, $T=1000$).}\label{tab:diff_path}
\centering\small
\begin{tabular}{lcccc}
\toprule
Step & LR/Schedule & Dataset & Epochs & FID $\downarrow$ \\
\midrule
Baseline & 2e-4, fixed & DEV & 50 & 100.2 \\
+ cosine LR & 2e-4, cosine & DEV & 100 & 78.3 \\
+ higher LR & 4e-4, cosine & DEV & 100 & 65.7 \\
\textbf{+ EMA (200ep)} & 4e-4, cosine & PROD & 200 & \textbf{29.0} \\
\bottomrule
\end{tabular}
\end{table}

% ---------------------------------------------------------------
\section{Qualitative Analysis and Discussion}\label{sec:final}

\begin{figure}[htbp]
\centering
\begin{subfigure}[b]{0.32\textwidth}
    \includegraphics[width=\textwidth]{figures/vae_samples.png}
    \caption{VAE ($\beta{=}0.05$)}
\end{subfigure}
\hfill
\begin{subfigure}[b]{0.32\textwidth}
    \includegraphics[width=\textwidth]{figures/dcgan_samples.png}
    \caption{DCGAN+SN}
\end{subfigure}
\hfill
\begin{subfigure}[b]{0.32\textwidth}
    \includegraphics[width=\textwidth]{figures/diff_samples.png}
    \caption{Diffusion+EMA}
\end{subfigure}
\caption{Random samples. VAE is blurry (\textbf{conditional mean optimization}); GAN is sharp; Diffusion has highest multi-scale richness.}
\label{fig:samples}
\end{figure}

\begin{table}[htbp]
\caption{Final PROD results (full dataset, 200 epochs, 10 seeds).}\label{tab:main}
\centering
\begin{tabular}{llcc}
\toprule
Model & Key configuration & FID $\downarrow$ & KID $\downarrow$ \\
\midrule
VAE & $\beta{=}0.05$, $d{=}128$ & 145.85 $\pm$ 1.27 & 0.1497 $\pm$ 0.009 \\
\textbf{DCGAN+SN} & ngf${=}128$, lat${=}100$ & \textbf{28.12 $\pm$ 0.59} & \textbf{0.0106 $\pm$ 0.003} \\
Diff.+EMA & $C{=}96$, DDIM & 28.97 $\pm$ 0.38 & 0.0144 $\pm$ 0.004 \\
\bottomrule
\end{tabular}
\end{table}

\noindent\textbf{Trade-off Analysis.}
The GAN is $\approx$\textbf{60$\times$ faster} than Diffusion for sampling. Both match at $32{\times}32$ because: (1) \textbf{resolution ceiling} limits diffusion's multi-scale advantage; (2) \textbf{metric bias} favors GANs which optimize inception-feature statistics directly.

\begin{figure}[htbp]
\centering
\begin{subfigure}[b]{0.32\textwidth}
    \includegraphics[width=\textwidth]{figures/vae_recon.png}
    \caption{VAE reconstructions}
\end{subfigure}
\hfill
\begin{subfigure}[b]{0.32\textwidth}
    \includegraphics[width=\textwidth]{figures/diff_recon.png}
    \caption{Diffusion reconstructions}
\end{subfigure}
\hfill
\begin{subfigure}[b]{0.32\textwidth}
    \includegraphics[width=\textwidth]{figures/diff_denoise.png}
    \caption{Denoising progress}
\end{subfigure}
\caption{Reconstructions and iterative denoising progress ($t{=}990 \to 0$).}
\label{fig:recon}
\end{figure}

% ---------------------------------------------------------------
\section{Conclusion}\label{sec:conclusion}
We conducted 100+ experiments. Findings: (1) $\beta \in [0.05, 0.1]$ avoids posterior collapse. (2) SN is essential for long GAN stability. (3) Cosine LR scheduling is model-family-dependent. (4) GAN and Diffusion achieve parity at $32{\times}32$ due to resolution ceilings.

% ---------------------------------------------------------------
\begin{credits}
\ackname{} Claude (Anthropic) assisted in code debugging and analysis.
\discintname{} No competing interests.
\end{credits}

\begin{thebibliography}{14}
\bibitem{binkowski2018} Binkowski, M. et al.: Demystifying MMD GANs. In: ICLR (2018)
\bibitem{dhariwal2021} Dhariwal, P., Nichol, A.: Diffusion beat GANs. In: NeurIPS (2021)
\bibitem{gulrajani2017} Gulrajani, I. et al.: Improved training of WGANs. In: NeurIPS (2017)
\bibitem{heusel2017} Heusel, M. et al.: TTUR for GANs. In: NeurIPS (2017)
\bibitem{higgins2017} Higgins, I. et al.: $\beta$-VAE. In: ICLR (2017)
\bibitem{ho2020} Ho, J. et al.: Denoising diffusion probabilistic models. In: NeurIPS (2020)
\bibitem{karras2018} Karras, T. et al.: StyleGAN architecture. arXiv:1812.04948 (2018)
\bibitem{kingma2013} Kingma, D.P., Welling, M.: Auto-encoding variational Bayes. (2013)
\bibitem{liao2022} Liao, P. et al.: The ArtBench dataset. arXiv:2206.11404 (2022)
\bibitem{miyato2018} Miyato, T. et al.: Spectral normalization for GANs. In: ICLR (2018)
\bibitem{oord2017} van den Oord, A. et al.: VQ-VAE. In: NeurIPS (2017)
\bibitem{radford2015} Radford, A. et al.: DCGAN. arXiv:1511.06434 (2015)
\bibitem{rombach2022} Rombach, R. et al.: Synthesis with LDMs. In: CVPR (2022)
\bibitem{song2020} Song, J. et al.: DDIM. arXiv:2010.02502 (2020)
\bibitem{extra} Kingma, D.P.: Adam: A Method for Stochastic Optimization. (2014)
\end{thebibliography}

\end{document}